% chapters/05/checklist.tex
\bagian{Checklist Pencapaian Kompetensi}

Bab ini menuntaskan \SUBCPMK{1.2}. Centang setiap butir berikut hanya apabila
Anda benar-benar telah melakukannya atau dapat menjelaskannya.

\begin{checklistbox}
\begin{checklist}
  \item {\bfseries (Indikator a)} Saya dapat menyatakan posisi dan kecepatan
        objek sebagai \code{pygame.Vector2}.
  \item Saya dapat menjelaskan mengapa gerak ke atas dinyatakan sebagai
        pengurangan nilai pada komponen tegak.
  \item Saya dapat membaca keadaan tombol dengan
        \code{pygame.key.get_pressed()} dan menyusunnya menjadi vektor arah.
  \item Saya dapat menjelaskan mengapa vektor arah perlu dinormalkan dan apa
        akibatnya apabila langkah itu dilupakan.
  \item Saya dapat menulis baris perhitungan perpindahan objek beserta
        penjelasan setiap unsurnya.
  \item {\bfseries (Indikator b)} Saya telah mengukur kecepatan objek pada 30
        dan 60 bingkai per detik dan mencatat hasilnya pada lembar kerja.
  \item Kecepatan terukur pada kedua laju bingkai itu bernilai sama pada
        program saya.
  \item Saya dapat menjelaskan mengapa jarak tempuh pada kedua pengukuran itu
        tidak sama persis walaupun kecepatannya sama.
  \item {\bfseries (Indikator c)} Saya dapat menyelaraskan posisi bertipe
        pecahan dengan \code{Rect} bertipe bilangan bulat.
  \item Saya dapat menjelaskan mengapa posisi tidak boleh disimpan sebagai
        bilangan bulat beserta akibatnya pada kecepatan rendah.
  \item Saya dapat menjelaskan perbedaan \code{round()} dan \code{int()} pada
        penyelarasan posisi.
  \item Program gerak saya berjalan pada dua laju bingkai berbeda dengan
        kecepatan objek yang sama.
  \item Catatan praktikum saya memuat tabel pengukuran kecepatan pada tiga
        laju bingkai.
\end{checklist}
\end{checklistbox}
